% DMRN+21: the official 2026 kit ships the 2025 style; its event header and
% footer are updated in the local copy, without changing its page layout.
\documentclass[a4paper]{article}
\usepackage[T1]{fontenc}
\usepackage{caption}
\usepackage{times,amsmath,graphicx,url,dmrn2026}
\usepackage{indentfirst}
\usepackage{xcolor}
\captionsetup[figure]{skip=2pt}\setlength{\intextsep}{3pt plus 1pt minus 1pt}
\setlength{\parskip}{0.1em}
\captionsetup[figure]{font=footnotesize,labelfont=footnotesize}
\captionsetup[table]{font=footnotesize,labelfont=footnotesize}
\setlength{\textfloatsep}{3pt plus 1pt minus 1pt}
\setlength{\belowcaptionskip}{0pt}
\renewcommand{\thesection}{\Roman{section}}
\newcommand{\head}[1]{\par\vspace{1pt}\noindent\textbf{#1.}\ \ignorespaces}
% Tighter display spacing. \ninept makes \normalsize an alias of \small, which
% resets the display skips at every size switch, so the reduced values are
% appended to \small itself.
\makeatletter
\g@addto@macro\small{\abovedisplayskip 4pt plus 1pt minus 1pt \belowdisplayskip 4pt plus 1pt minus 1pt
 \abovedisplayshortskip 2pt \belowdisplayshortskip 2pt plus 1pt}
\makeatother
% No extra space between reference entries.
\let\oldthebibliography\thebibliography
\renewcommand{\thebibliography}[1]{\oldthebibliography{#1}\setlength{\itemsep}{0pt}\setlength{\parskip}{0pt}}
% Layout economies after the DMRN+15 auraloss one-pager: tighter title block,
% 7 pt footer close to the page edge, taller text block. The style file is unchanged.
\setlength{\textheight}{266truemm}\setlength{\footskip}{0.2in}
\cfoot{\fontsize{7}{8}\selectfont\textsc{DMRN+21: DIGITAL MUSIC RESEARCH NETWORK ONE-DAY WORKSHOP 2026, KING'S COLLEGE LONDON, FRI 11 DECEMBER 2026}}
\makeatletter
\def\@maketitle{\newpage
 \null
 \begin{minipage}{0.5\textwidth}
 \textsc{DMRN+21: Digital Music Research Network\\One-day Workshop 2026}\vspace{3mm}\\\textsc{King's College London}\\\textsc{Fri 11 December 2026}
 \end{minipage}%
 \begin{minipage}{0.5\textwidth}
 \hfill\hbox{\includegraphics[width=6cm]{DMRN_Logo_high_res.jpg}}
 \end{minipage}
 \vskip 0.4em \begin{center}
 {\large \bf \@title \par} \vskip 0.35em {\large
\begin{tabular}[t]{c}\@name \\ \@address
 \end{tabular}
 \par}
 \end{center}
 \par
 \vskip 0.2em}
\makeatother

\title{Efficient Backpropagation for\\Finite-Difference Nonlinear String Synthesis}
\name{\normalfont Pablo Tablas de Paula, Emmanouil Benetos and Joshua D. Reiss\thanks{This work was supported by UKRI grant EP/S022694.}}
\address{\small Centre for Digital Music, Queen Mary University of London, p.tablasdepaula@qmul.ac.uk}

\begin{document}
\ninept
\maketitle
\begin{sloppy}
\begin{abstract}
Physical models of strings solved with finite-difference time-domain (FDTD)
schemes are rarely used in differentiable sound synthesis, because
backpropagating through tens of thousands of time steps is slow and
memory-hungry. We show that the adjoint method with checkpointing, long
established in physics simulation, removes this obstacle for a nonlinear
string whose tension rises as it vibrates (the Kirchhoff--Carrier model),
once scheme and adjoint are compiled. For a batch of 32 one-second strings, one
sound-matching gradient step takes 3.3 s with an off-the-shelf compiler
(JAX) and 0.38 s in a hand-written GPU kernel, against 263 s for uncompiled
automatic differentiation with the same checkpoints. Gradient checks and
applications are in the online
supplement.\footnote{\fontsize{6}{7}\selectfont\url{ptablasdpaula.github.io/DMRN21-Kirchhoff_Carrier_String}}
\end{abstract}

\section{Introduction}
Finite-difference time-domain (FDTD) string synthesis is increasingly
leaner~\cite{bilbao2024guitar} and runs on the GPU~\cite{lee2024string}, yet differentiable
frameworks~\cite{hayes2024review} largely avoid it because
backpropagating through a time-stepping scheme builds deep and expensive
recurrent graphs. Hence, the literature has favoured
digital waveguides~\cite{tablas2026sound} and modal
synthesis~\cite{diaz2025}. Physics simulation has long met the same obstacle
with the adjoint method and checkpointing, now the default of
differential-equation solvers such as Diffrax~\cite{kidger2021}, and
hand-written adjoints already serve recursive audio
filters~\cite{yu2024allpole}. We apply this recipe to a nonlinear FDTD
string and measure what each ingredient buys.

\section{Model and method}
We model the Kirchhoff--Carrier string~\cite[Sec.~8.1]{bilbao2009}, a stiff
and lossy string whose tension rises with its stretch, simply supported at
both ends~\cite{russo2026}:
\begin{equation}
\begin{aligned}
\rho A\frac{\partial^2u}{\partial t^2}={}&
 \bigg[T_0+\frac{EA}{2L}\!\int_0^L\!\Big(\frac{\partial u}{\partial x}\Big)^{\!2}dx\bigg]
 \frac{\partial^2u}{\partial x^2}
 -EI\frac{\partial^4u}{\partial x^4}\\
 &-2\rho A\Big(\sigma_0-\sigma_1\frac{\partial^2}{\partial x^2}\Big)
 \frac{\partial u}{\partial t},
\end{aligned}
\label{eq:pde}
\end{equation}
with displacement $u(x,t)$, density $\rho$, cross-section $A$, Young's
modulus $E$, second moment of area $I$, length $L$ and damping
$\sigma_0,\sigma_1$. The bracket is the tension: its rest value $T_0$ plus a
term that grows with the string's stretch, one number for the whole string,
which is what bends the pitch of a hard pluck.

\head{Time stepping}
We integrate (\ref{eq:pde}) with the energy-conserving scheme
of~\cite[Sec.~8.1.2]{bilbao2009}, extended with stiffness and loss
in~\cite{russo2026}. Because the nonlinearity is a single number, each
step solves for it in closed form, then advances every spatial sine mode
independently.

\head{Checkpointed adjoint}
Backpropagation visits the time steps in reverse and needs the state at
each one. Automatic differentiation (autodiff) stores them all, so its
memory grows with duration. With checkpointing, the state is stored only
every 256 steps: the reverse pass re-runs one segment at a time from its
checkpoint, then steps the adjoint of the scheme, derived by hand, back
through it. Memory is the checkpoints plus one segment, for the price of a
second forward pass, and the gradient is still that of the discrete scheme.
Both passes run in a single GPU (CUDA) kernel, in parallel over modes and
batch items.

\head{Spatial stencil}
To reduce numerical dispersion, we use a fourth-order five-point
stencil~\cite[Sec.~6.3.1]{bilbao2009} for the spatial differences. For the benchmark pluck,
waveform error against a converged solution from~\cite{russo2026} falls from
11.7\% to 1\%, even with oversampling reduced from 3$\times$ to 2$\times$.

\section{Evaluation}
Each evaluation is one sound-matching step: a batch of strings is simulated
for 1 s, the mean absolute error of each waveform against the target's is
computed, and backpropagated to ten parameters:
$\rho$, $T_0$, $r$ ($A=\pi r^2$, $I=\pi r^4/4$), $E$, $L$, $\sigma_0$,
$\sigma_1$, the pluck amplitude and the pluck and pickup positions. The target
is a steel string ($\rho=8050$ kg/m$^3$, $T_0=75$ N, $r=0.356$ mm, $E=174$ GPa,
$L=1$ m, $\sigma_0=0.92$ s$^{-1}$, $\sigma_1=2.9\times10^{-4}$ m$^2$/s), plucked
at its centre with a 2 mm raised cosine of half-width $0.24L$, its
displacement read at its midpoint; the batch perturbs every parameter by up to
$\pm3\%$. The scheme uses 176 intervals at 88.2 kHz, decimated to 44.1 kHz,
on an A100 80 GB in float64. Fig.~\ref{fig:benchmark} compares four ways of
obtaining the same gradient: \emph{Autodiff}, eager PyTorch autodiff of the
scheme on the grid; \emph{Autodiff, checkpointed}, the same with 256-step
checkpoints; \emph{Compiled autodiff}, the checkpointed scheme in the sine
basis, compiled by JAX; and the \emph{Adjoint kernel} above. Checkpointing
fixes memory: at batch 32, autodiff falls from 44.9 GB to 0.9 GB, for 1.5
times the time. Compilation fixes time: at the same checkpoints, PyTorch's
compiler, limited to 16-step blocks, is ${\sim}10$ times faster than eager
PyTorch; JAX, compiling the whole loop, ${\sim}100$ times; and our kernel a
further 5 to 9 times.

\begin{figure}[t]
 \centering
 \includegraphics[width=\columnwidth]{figures/benchmark-dots-label.pdf}
 \caption{Peak memory and time per step at the first and last batch sizes.}
 \label{fig:benchmark}
\end{figure}

{\footnotesize
\begin{thebibliography}{10}
\bibitem{bilbao2024guitar} S.~Bilbao \emph{et~al.}, ``Real-time guitar synthesis,'' \emph{Proc. DAFx}, 2024.
\bibitem{lee2024string} J.~W.~Lee \emph{et~al.}, ``String sound synthesizer on GPU-accelerated finite difference scheme,'' \emph{Proc. ICASSP}, 2024.
\bibitem{hayes2024review} B.~Hayes \emph{et~al.}, ``A review of differentiable digital signal processing for music and speech synthesis,'' \emph{Front. Sig. Proc.}, v.~3, 1284100, 2024.
\bibitem{tablas2026sound} P.~Tablas~de~Paula \emph{et~al.}, ``Sound matching with a differentiable Karplus-Strong algorithm,'' \emph{Proc. DAFx}, 2026.
\bibitem{diaz2025} R.~Diaz and M.~Sandler, ``Fast differentiable modal simulation of non-linear strings, membranes, and plates,'' \emph{Proc. DAFx}, 2025.
\bibitem{kidger2021} P.~Kidger, ``On neural differential equations,'' DPhil thesis, Oxford, 2021.
\bibitem{yu2024allpole} C.-Y.~Yu \emph{et~al.}, ``Differentiable all-pole filters for time-varying audio systems,'' \emph{Proc. DAFx}, 2024.
\bibitem{bilbao2009} S.~Bilbao, \emph{Numerical Sound Synthesis: Finite Difference Schemes and Simulation in Musical Acoustics}, Wiley, 2009.
\bibitem{russo2026} R.~Russo \emph{et~al.}, ``Numerical convergence of the scalar auxiliary variable method applied to nonlinear stiff string models,'' \emph{Nonlinear Dyn.}, v.~114, 2026.
\end{thebibliography}}
\end{sloppy}
\end{document}
